\documentclass[11pt,a4paper]{article}
\usepackage[utf8]{inputenc}
\usepackage[T1]{fontenc}
\usepackage{geometry}
\usepackage{graphicx}
\usepackage{booktabs}
\usepackage{caption}
\usepackage{array}
\usepackage{amsmath,amssymb,amsthm}
\usepackage[numbers]{natbib}
\usepackage{xurl}
\usepackage[colorlinks=true,linkcolor=blue,citecolor=blue,urlcolor=blue]{hyperref}
\usepackage{listings}

\geometry{
  left=2.5cm,
  right=2.5cm,
  top=2.5cm,
  bottom=2.5cm
}

\newcolumntype{L}[1]{>{\raggedright\arraybackslash}p{#1}}
\emergencystretch=2em

\title{\textbf{NRR-Patterns: Condition-Bounded Comparison of Five Delayed-Commitment Patterns in Stateful LLM Systems}}
\author{
  Kei Saito\thanks{ORCID: 0009-0006-4715-9176} \\
  Independent Researcher, Japan \\
  \texttt{kei.saito.research@gmail.com}
}
\date{April 2026 \\[0.5em]
\textit{Part of the Non-Resolution Reasoning (NRR) research program.}}

\begin{document}
\maketitle

\begin{center}
\textcopyright\ 2026 Kei Saito.
Licensed under CC BY 4.0.\\
\url{https://creativecommons.org/licenses/by/4.0/}
\end{center}

\begin{abstract}
Delayed-commitment recommendations in large language model systems are often discussed too broadly or reused as if they transfer uniformly across providers and task conditions. This paper compares five delayed-commitment patterns on a fixed two-run matched paired comparison spine with live provider APIs and then reports that released output surface through a provider-separated stability read, a separated temperature-robustness member, and a selected Stage B boundary-honesty carry-forward. The main comparison map is carried by the fixed matched-run spine, while sign-flip and non-effective regions are carried by separately labeled selected Stage B artifacts rather than folded into the headline comparison itself. On the fixed matched-run spine, Hierarchical Resolution, Delta Update, Lazy Expansion, Conditional Refinement, and Branch Preservation are all positive relative to their own paired baselines at the provider-balanced aggregate level, while Branch remains only weakly positive on that aggregate read. The provider-separated view is less uniform: Branch already contains negative Anthropic members on the fixed spine before failing Stage2, and Conditional weakens under Stage2 while retaining a positive provider-balanced aggregate despite an Anthropic-side reversal. The separated temperature-robustness member then shows that Branch is the clearest fragile case, but the failure shape is tail-risk-dominated rather than provider-wide: the Gemini-side negative mean is driven mainly by one catastrophic row, while the remaining Gemini rows are near-neutral in aggregate. A compact compression comparison over a retained three-pattern subset shows that the observed gains for that tested subset are not reducible to generic prompt shortening: principle-only remains clearly positive within each tested family-relative contrast, whereas compression-only weakens or fails for key cases. Selected Stage B boundary checks then show that a single aggregate comparison view can hide provider-sensitive reversals and non-effective cells. On the retained key metrics, Conditional has the clearest positive provider-mean carry-forward, but that carry-forward remains sparse rather than uniformly active, while Lazy is the clearest degraded case and Hierarchical, Delta, and Branch show provider-sensitive weakening or reversal cells. A separate selected non-effective slice then shows all-provider-zero cells for Conditional, Delta, and Branch on internal-coherence and trace-completeness diagnostics. The reviewer-facing package audits this fixed output surface rather than re-exporting a full rerun workspace. The resulting contribution is a condition-bounded comparison whose comparison headline is kept honest by explicit provider-separated stability, tail-risk disclosure, robustness, and separately labeled sign-flip and non-effective reporting rather than by a same-scale cross-family ranking or universal deployment rule.

\end{abstract}

\noindent\textbf{Keywords:} Non-Resolution Reasoning, delayed commitment, stateful LLM systems, paired comparison, provider-separated stability, boundary honesty

\section{Introduction}
\label{sec:intro}

Delayed commitment is useful in many stateful LLM workflows, but recommendations built from that idea are easy to overgeneralize. A pattern that reduces cost in one setting may weaken, reverse, or become non-effective under another provider or operating condition. The main problem for this paper is therefore not whether delayed commitment can ever help, but how to compare five tested delayed-commitment patterns on one auditable fixed matched-run comparison spine while keeping the comparison evidence separate from the selected Stage B honesty carry-forward used for sign-flip and non-effective reporting. In this paper, the headline comparison comes from the fixed matched \(t=0.0\) experiment spine and the released outputs traced from it; the reviewer-facing ZIP is an audit surface for those outputs rather than an independent rerun handoff, and the selected Stage B honesty layer remains a separate carry-forward caveat surface rather than part of the headline spine.

This NRR-Patterns manuscript is the comparison-led pattern paper in the current NRR spine. Fixed public lines in that spine establish first principles and text-to-state mapping~\cite{saito2025nrr,saito2026phi}. Additional IME, Transfer, and Coupled lines remain series-path references rather than fixed evidentiary supports for the claims made here~\cite{saito2026ime,saito2026transfer,saito2026coupled}. The role here is narrower and more practical: compare five delayed-commitment patterns on a fixed matched-run paired spine, then report that audited output surface through provider-separated stability reads, a separated temperature-robustness member, and a separately labeled selected Stage B honesty layer. In the downstream series read, this paper is the upstream comparison-and-boundary-honesty layer for later calibration and closure work, not a universal endpoint.

NRR is not an anti-LLM framework.
NRR does not replace standard LLM use.
NRR optimizes when to commit and when to defer, under explicit conditions.

\textbf{Series Consistency Statement:}
Across the NRR series, unsupported interpretations are not injected at a fixed evidence state, hypothesis-set expansion is evidence-gated, and unobserved possibilities remain explicitly unresolved. In NRR-Patterns, this policy is operationalized by keeping the main five-pattern comparison separate from selected boundary-honesty carry-forward, and by reporting unstable, sign-flip, and non-effective regions explicitly rather than hiding them behind winner-only summaries.

\textbf{Series Extensibility Statement:}
NRR should be read here as a derivation-and-evaluation framework, not as a closed recipe. The five delayed-commitment patterns in NRR-Patterns are an evaluated set under the current fixed matched-run paired spine plus explicitly separated robustness and honesty layers, not an exhaustive set of admissible patterns. Additional patterns or operator-family extensions are in scope only when they are tested under the same paired-unit discipline and reported with equally explicit condition-bounded boundary language.

NRR-Patterns compares five delayed-commitment patterns on a fixed matched-run paired spine and then adds separate provider-separated stability, temperature-robustness, and selected Stage B honesty layers to show where that map weakens, reverses, or becomes non-effective.

The five canonical patterns are:
\begin{enumerate}
  \item Hierarchical Resolution
  \item Lazy Expansion
  \item Branch Preservation
  \item Delta Update
  \item Conditional Refinement
\end{enumerate}

These are treated here as delayed-commitment \emph{patterns}, not as the default meaning of lower-level \emph{operator} language used elsewhere in the series. That distinction matters because other papers in the series path also use \((\sigma/\delta,\text{target})\)-style operator language for interface-level update primitives, but those mentions are contextual here rather than evidentiary supports for the present pattern-comparison claims~\cite{saito2026phi,saito2026transfer,saito2026coupled}. The current paper instead evaluates pattern-level delayed-commitment designs.

\subsection{Contributions}

\begin{enumerate}
  \item A fixed matched-run comparison of five delayed-commitment patterns on two matched \(t=0.0\) runs, released here as an auditable output surface rather than a full rerun bundle.
  \item A separated provider-stability and temperature-robustness read showing where the comparison headline preserves, weakens, or breaks under provider-local condition shifts.
  \item A condition-bounded pattern map whose positive and unstable comparison regions are carried by the fixed matched-run spine rather than by a same-scale cross-family ranking.
  \item A separately labeled selected Stage B honesty layer for sign-flip and non-effective regions, showing where aggregate comparison reads can hide reversals or weak cells.
  \item A compression control over a retained three-pattern subset showing that those tested gains are not reducible to generic prompt shortening.
\end{enumerate}

\section{Related Work}
\label{sec:related}

\textbf{Delayed commitment and stateful interaction.}
Prior NRR work covers the failure mode and text-to-state mapping needed for stateful execution on stateless APIs in fixed public lines~\cite{saito2025nrr,saito2026phi}. Additional IME, Transfer, and Coupled references are retained here only to mark the broader series path and not as fixed evidentiary supports for the present comparison claims~\cite{saito2026ime,saito2026transfer,saito2026coupled}. The present paper does not reopen those layers. Instead, it compares multiple delayed-commitment patterns built on top of them and asks how that comparison should be read under explicit provider-sensitive boundaries.

\textbf{Efficiency and compression.}
Context-efficiency work, including retrieval compression and prompt-side compaction, targets token cost and context footprint directly~\cite{xu2023recomp,chevalier2023compress}. This paper overlaps in efficiency interest but differs in object: the main question is not simply whether prompts can be shortened, but whether delaying specific commitments yields a usable pattern map that remains condition-bounded once provider-sensitive boundary checks are surfaced. That is why the compression comparison is kept as a control rather than as the paper's main identity.

\textbf{Boundary-aware evaluation.}
Evaluation-framework work such as HELM, MT-Bench, and AlpacaEval shows that reported quality or preference conclusions can shift with benchmark protocol, judge setup, and reporting frame~\cite{liang2022helm,zheng2023mtbench,dubois2024alpacaeval}.

\textbf{Prompting and calibration sensitivity.}
Prompting and calibration work likewise shows that effect direction can change with reasoning prompts, aggregation strategy, and label-bias controls~\cite{kojima2022zeroshot,wang2023selfconsistency,zhao2021calibrate}. This manuscript adopts the shared lesson in a narrower form: it keeps one exact-replication paired spine, then reports a separated temperature-robustness member plus selected provider-sensitive boundary evidence so the comparison map is not over-read as a same-scale cross-family ranking.

\section{Pattern Set and Evaluation Logic}
\label{sec:method}

\subsection{Pattern Set and Paired Baselines}

Each pattern is evaluated under paired baseline-vs-pattern comparison with fixed scenario definitions and prompt templates inherited from the current released evidence surfaces. Table~\ref{tab:pattern_catalog} states the pattern set, its baseline contrast, and the current evidential role in this manuscript.

\begin{table}[t]
\centering
\small
\begin{tabular}{L{3.0cm}L{4.2cm}L{2.5cm}L{4.0cm}}
\toprule
\textbf{Pattern} & \textbf{Paired baseline contrast} & \textbf{Delayed-commitment axis} & \textbf{Current paper role} \\
\midrule
Hierarchical Resolution & flat full-option selection vs category-then-tool selection & decision granularity & positive family-relative benchmark with large within-family reduction \\
Lazy Expansion & eager full retrieval vs summary-first selective retrieval & content expansion & positive family-relative benchmark with later boundary sensitivity \\
Branch Preservation & no branch memory vs checkpointed branch restoration & branch destruction & provider-local instability case on the comparison surface \\
Delta Update & full-state resend vs changed-item-only updates & state retransmission & positive family-relative benchmark \\
Conditional Refinement & detailed-by-default vs coarse-first refinement & output detail & positive family-relative benchmark with sparse boundary carry-forward \\
\bottomrule
\end{tabular}
\caption{Canonical pattern catalog for NRR-Patterns. The comparison-led manuscript centers these five patterns rather than treating the paper as a broad standalone foundations manuscript.}
\label{tab:pattern_catalog}
\end{table}

\subsection{Primary and Supporting Metrics}

The headline comparison metric is \emph{provider-balanced weighted reduction} on the exact-replication spine. The paper-local comparison contract for this quantity is frozen in \path{repro/paper7_split_comparison_contract_v12_2026-04-01.md}. For a given pattern and carried-forward run, each provider first contributes its own weighted reduction percentage, \(100 \times \sum_i (b_i - p_i) / \sum_i b_i\), over the paired rows in that run's shipped robust source member, where \(b_i\) and \(p_i\) denote the baseline and pattern token totals for the same paired row. The run-level quantity \path{weighted_percent_provider_balanced} is then the simple arithmetic mean of those provider-level weighted percentages across the shipped providers for that run member. Table~\ref{tab:comparison_summary} reports the two exact-replication run values, their exact-replication mean, and the separated Stage2 temperature-robustness member, while Table~\ref{tab:provider_separated_stability} reports the compact provider-separated stability summary derived from the same paired rows. Supporting comparison metrics remain weighted global reduction, median reduction, trimmed reduction, and cost-per-success delta. These supporting metrics remain useful, but they do not replace the headline role of the exact-replication provider-balanced read in the main comparison section.

\subsection{Boundary-Honesty Metrics}

Selected carried-forward boundary metrics are misinterpretation improvement and rework improvement. Here ``Stage B'' means the historical all-pattern confirmatory boundary-resolution suite carried in the `nrr-boundary' package, from which this manuscript retains only a small honesty-facing subset rather than the full program. The paper-local selection rule for carrying Stage B into NRR-Patterns is frozen in \path{repro/paper7_stageb_selection_rule_v1_2026-03-31.md}. This is an integration rule for the comparison-led manuscript rather than a claim that this manuscript originally preregistered the full Stage B program. The retained pair is the smallest downstream quality/cost pair that can make the five-pattern comparison honest without turning NRR-Patterns into a second boundary-led manuscript. Their role is not to replace the main token-efficiency center, but to show where the pattern map becomes provider-sensitive, unstable, or non-effective. Additional diagnostics such as trace completeness, update consistency, internal coherence error rate, self-contradiction rate, and selected stability diagnostics remain available as secondary/package-facing surfaces when they materially sharpen the honesty of the paper-level read.

\subsection{Reporting Rules and Validity Limits}

Use provider-balanced results for the headline comparison and provider-separated results for boundary honesty where aggregation could hide reversals.

The canonical result-language ladder is:
\begin{itemize}
  \item positive region
  \item unstable region
  \item sign-flip region
  \item non-effective region
\end{itemize}

The main comparison surface primarily supports the positive/unstable read. Sign-flip and non-effective labels are carried by selected boundary checks where those regions are directly observed. Accordingly, these labels should not be misread as if they were all outputs of the headline comparison metric alone.

\section{Experimental Surfaces}
\label{sec:surfaces}

\subsection{Main Comparison Surface}

The primary experimental spine uses the exact-replication five-pattern comparison surface currently carried by the released NRR-Patterns line. Here, exact replication continues to mean the fixed auditable matched-run output surface named in the shipped artifacts, not a rerun-capable environment. The shipped provider/model scope is Anthropic Claude Sonnet 4 (\texttt{claude-sonnet-4-20250514}), Gemini 2.0 Flash (\texttt{gemini-2.0-flash}), and OpenAI GPT-4o mini (\texttt{gpt-4o-mini-2024-07-18}):
\begin{itemize}
  \item Stage1 main (\(t=0.0\))
  \item Stage1 rep2 (\(t=0.0\))
\end{itemize}

These two \(t=0.0\) runs define the headline exact-replication spine for the integrated comparison read. The integrated source-of-truth name freeze for these comparison surfaces is recorded in \path{repro/paper7_integrated_comparison_claims_manifest_v6_2026-04-01.csv}, and the paper-local fairness/aggregation contract for separating headline replication from robustness is frozen in \path{repro/paper7_split_comparison_contract_v12_2026-04-01.md}. The current manuscript-facing comparison headline summary is \path{results/analysis/paper7_integrated_comparison_headline_summary_v2_2026-03-31.csv}, the compact provider-separated stability summary is \path{results/analysis/paper7_provider_separated_stability_summary_v2_2026-04-01.csv}, and the headline source members are:
\begin{itemize}
  \item \path{results/paper7_results_allproviders_clean.zip::results/analysis/paper7_allproviders_pattern_summary_robust_t00.csv}
  \item \path{results/paper7_results_stage1_t00_rep2.zip::results/analysis/paper7_allproviders_pattern_summary_robust_t00.csv}
\end{itemize}
The carried-forward temperature-robustness member is then:
\begin{itemize}
  \item \path{results/paper7_results_stage2_t03.zip::results/analysis/paper7_allproviders_pattern_summary_robust_t03.csv}
\end{itemize}
The manuscript therefore exposes one exact-replication headline summary for inspection while preserving the robustness member for drill-down rather than collapsing all three members into one undifferentiated mean.

\subsection{Compression Comparison}

The compression comparison remains in the main manuscript as a compact interpretive control. Its role is to answer the reviewer question: are the observed gains just generic prompt compression? The canonical integrated name for this control is also frozen in \path{repro/paper7_integrated_comparison_claims_manifest_v6_2026-04-01.csv}. The current aggregated control surface is \path{results/analysis/paper7_comparison_summary_r2_aggregated.csv}, with per-provider drill-down in \path{results/analysis/paper7_comparison_summary_r2_by_provider.csv} and raw comparison traces inside the three \path{results/paper7_compare_t00_r2_anthropic.zip}, \path{results/paper7_compare_t00_r2_openai.zip}, and \path{results/paper7_compare_t00_r2_gemini.zip} artifacts.

\subsection{Selected Stage B Boundary-Resolution Surface}

Selected Stage B evidence enters this manuscript as a carried-forward boundary-honesty layer, not as the paper's primary confirmatory center. In this manuscript, ``Stage B'' refers to the historical `nrr-boundary' confirmatory boundary-resolution suite over the same five delayed-commitment families, but this manuscript carries forward only the selected honesty-facing subset named here. The integrated source-of-truth name freeze for this layer is recorded in \path{repro/paper7_integrated_boundary_claims_manifest_v1_2026-03-29.csv}, and the paper-local rule for why these artifacts are retained here is frozen in \path{repro/paper7_stageb_selection_rule_v1_2026-03-31.md}. That rule is a paper-local integration rule rather than a manuscript preregistration freeze, and it retains \texttt{misinterpretation\_rate} plus \texttt{rework\_cost\_tokens} as the smallest downstream quality/cost pair needed to make the comparison honest. In the integrated package, the carried-forward Stage B read is traceable to:
\begin{itemize}
  \item \path{nrr-boundary/stats/stageb_all/mst_provider_balanced.csv}
  \item \path{nrr-boundary/stats/stageb_all/mst_provider_separated.csv}
  \item \path{nrr-boundary/stats/stageb_all/mst_sign_flip_boundaries.csv}
  \item \path{nrr-boundary/stats/stageb_all/stageb_confirmatory_tests_holm30.csv}
\end{itemize}
These are used here only to make the comparison map honest about provider-sensitive weakening, reversals, and non-effective cells. They are not used to replace the token-efficiency headline or to turn this manuscript into a second boundary-led manuscript.

\subsection{Secondary / Package-First Surfaces}

Ordered-combination follow-up, rep1-vs-rep3 direction details, and no-position ablation are not part of the default mainline. Ordered composition remains secondary because it introduces a separate provider-by-order interaction story that can easily dominate the paper if over-promoted. The no-position ablation remains useful as a robustness check, but it is not required for the main comparison claim hierarchy. In the current integrated read, both therefore remain appendix/package-first by default.

\section{Results: Pattern Comparison}
\label{sec:results_comparison}

\subsection{Provider-Balanced Pattern Map}

Across the two exact-replication comparison runs, provider-balanced weighted reduction remains strongly heterogeneous rather than uniform. Table~\ref{tab:comparison_summary} gives the current manuscript-facing summary from the integrated comparison headline summary named in Sec.~\ref{sec:surfaces}, while keeping the Stage2 temperature-robustness member separate from the headline mean.
\begin{table}[t]
\centering
\small
\begin{tabular}{L{2.2cm}L{1.3cm}rrrr}
\toprule
\textbf{Pattern} & \textbf{Rows} & \textbf{S1 main} & \textbf{S1 rep2} & \textbf{Rep. mean} & \textbf{S2} \\
\midrule
Hierarchical & 45 & \(+56.752\%\) & \(+57.606\%\) & \(+57.179\%\) & \(+55.737\%\) \\
Delta & 18 & \(+39.074\%\) & \(+38.956\%\) & \(+39.015\%\) & \(+38.980\%\) \\
Lazy & 45 & \(+37.951\%\) & \(+37.459\%\) & \(+37.705\%\) & \(+38.066\%\) \\
Conditional & 18 & \(+37.011\%\) & \(+35.353\%\) & \(+36.182\%\) & \(+32.235\%\) \\
Branch & 45 & \(+2.452\%\) & \(+3.081\%\) & \(+2.766\%\) & \(-21.713\%\) \\
\bottomrule
\end{tabular}
\caption{Provider-balanced weighted reduction under the split comparison contract. The headline mean uses only the two exact-replication \(t=0.0\) runs, while Stage2 \((t=0.3)\) is reported separately as the robustness member. The surface is not evidence-symmetric across patterns: Hierarchical, Lazy, and Branch carry \(n_{\mathrm{all}}=45\) paired rows per run, whereas Delta and Conditional carry \(n_{\mathrm{all}}=18\). The magnitudes are not directly comparable across families as a same-scale ranking because each family uses its own paired baseline contrast and token-cost floor.}
\label{tab:comparison_summary}
\end{table}

The resulting exact-replication read should therefore be interpreted with explicit paired-row-volume asymmetry, non-comparable family baseline contrasts, and provider-separated stability in view:
\begin{itemize}
  \item Hierarchical, Delta, and Lazy remain positive for Anthropic, Gemini, and OpenAI across the two exact-replication runs and the separated robustness member.
  \item Conditional is positive for all three providers on exact replication, but the Stage2 robustness member contains an Anthropic-side reversal even though the provider-balanced aggregate stays positive.
  \item Branch is aggregate-positive on exact replication but already contains negative Anthropic exact-replication members, then fails the separated robustness member through a Gemini-led, tail-risk-dominated collapse.
\end{itemize}

\subsection{Replication and Temperature Robustness}

The integrated paper retains the Stage2 member because the main issue is not just exact replication but also whether the same paired comparison object survives a controlled condition shift. That robustness member must be read separately rather than folded into the headline mean. Hierarchical, Delta, and Lazy remain positive for all three providers in both exact-replication runs and in Stage2. Conditional, by contrast, weakens at \(t=0.3\) and flips negative on Anthropic even though the provider-balanced Stage2 aggregate remains positive. Branch is the clearest failure case: it is slightly aggregate-positive in both \(t=0.0\) runs despite negative Anthropic exact-replication members, then turns sharply negative in Stage2 because Gemini contains a tail-risk event that pulls down the provider-balanced aggregate rather than producing a uniform all-provider collapse. That is why the integrated paper reads Branch as a fragile provider-local case rather than as a robust fifth positive family benchmark.

\begin{figure}[t]
\centering
\includegraphics[width=0.92\textwidth]{../figures/fig_patterns_weighted_3runs.png}
\caption{Provider-balanced weighted reduction across the two exact-replication runs plus the separated Stage2 robustness member. Error bars show the observed range over these members for each pattern. The plotted magnitudes remain family-relative against different paired baselines and should not be read as a same-scale cross-family ranking.}
\label{fig:paper7_comparison_3runs}
\end{figure}

\subsection{Provider-Separated Stability Read}

The split comparison surface already shows that exact-replication positivity is not the same as universal safety. Table~\ref{tab:provider_separated_stability} makes that point visible inside the main text rather than leaving it only in the shipped ZIP members. Branch is the clearest example: it stays slightly aggregate-positive on the exact-replication spine, but Anthropic is already negative in both exact-replication runs, and the Stage2 negative turn is driven by Gemini rather than by uniform all-provider collapse. Conditional is less severe but still illustrates the same principle by weakening materially at \(t=0.3\) while carrying an Anthropic-side reversal inside an aggregate-positive provider-balanced read. The other three patterns are substantially more stable across the split comparison read, but even there the correct interpretation is still condition-bounded rather than universal. This is why NRR-Patterns stops at a comparison-led family benchmark map and then carries boundary honesty separately, rather than converting the comparison table into a deployment-wide ranking.

\begin{table}[t]
\centering
\small
\begin{tabular}{L{1.85cm}L{2.7cm}L{2.7cm}L{2.7cm}L{3.75cm}}
\toprule
\textbf{Pattern} & \textbf{S1 main (A/G/O)} & \textbf{S1 rep2 (A/G/O)} & \textbf{S2 (A/G/O)} & \textbf{Comparison-side read} \\
\midrule
Hierarchical & A +41.3 / G +71.0 / O +57.9 & A +41.3 / G +74.2 / O +57.3 & A +41.7 / G +68.2 / O +57.4 & all-provider positive across the exact-replication spine and the separated robustness member \\
Delta & A +33.1 / G +35.9 / O +48.2 & A +33.7 / G +34.6 / O +48.6 & A +32.5 / G +38.9 / O +45.6 & all-provider positive across the comparison surface, but with lower support volume than Hierarchical/Lazy/Branch \\
Lazy & A +47.6 / G +41.1 / O +25.1 & A +47.6 / G +40.7 / O +24.1 & A +47.9 / G +42.1 / O +24.2 & all-provider positive on comparison and robustness, with later degradation appearing in the boundary layer \\
Conditional & A +6.8 / G +64.1 / O +40.1 & A +6.3 / G +61.0 / O +38.8 & A -4.1 / G +58.0 / O +42.8 & exact-replication all-provider positive, but Stage2 contains an Anthropic-side reversal inside an aggregate-positive read \\
Branch & A -0.1 / G +6.2 / O +1.2 & A -1.1 / G +6.4 / O +4.0 & A +5.3 / G -74.3 / O +3.9 & aggregate-positive exact replication with Anthropic instability, then Gemini-led Stage2 failure \\
\bottomrule
\end{tabular}
\caption{Compact provider-separated weighted reduction summary, derived from the reviewer-facing provider-separated stability summary artifact in the shipped audit surface. Provider order is Anthropic/Gemini/OpenAI. This table makes the comparison-side instability visible in the main text instead of leaving it only inside the shipped paired-delta ZIP members.}
\label{tab:provider_separated_stability}
\end{table}

The Branch failure shape is also narrower than a broad provider-wide collapse. In the Stage2 Gemini paired-delta member \path{results/paper7_results_stage2_t03.zip::results/analysis/paper7_prod_gemini_clean_t03_paired_deltas.csv}, the worst Branch row is \texttt{scenario3 / repetition3}, where the token count changes from \texttt{3372} to \texttt{38998} (\(-1056.524318\%\)). The Gemini Branch provider-weighted Stage2 value is \(-74.326453\%\) with that row included, but the compact provider-separated summary also records that it would be only \(-0.530303\%\) without that one worst row. The appropriate reading is therefore tail-risk-dominated robustness failure, not a broad all-row or all-provider collapse.

\subsection{Pattern Status Ladder}

The end of the comparison section collapses the mainline read into one compact status ladder. Table~\ref{tab:status_ladder} states the current integrated-paper summary.

\begin{table}[t]
\centering
\small
\begin{tabular}{L{2.2cm}L{1.2cm}L{2.6cm}L{5.5cm}}
\toprule
\textbf{Pattern} & \textbf{Rows} & \textbf{Split status} & \textbf{Boundary note} \\
\midrule
Hierarchical & 45 & robust positive region & remains strong on the exact-replication spine, but later boundary views still matter for provider-sensitive weakening cells \\
Lazy & 45 & robust positive region & exact-replication positive and robustness-positive, but carried-forward boundary views expose weaker and degrading cells \\
Branch & 45 & aggregate-positive instability / tail-risk robustness failure & exact replication is only aggregate-positive because Anthropic is already negative, and the separated robustness member fails through a Gemini-led tail-risk event \\
Delta & 18 & robust positive region & exact-replication positive and robustness-positive, but selected boundary checks expose weaker cells \\
Conditional & 18 & aggregate-positive with Stage2 provider reversal & clearest positive provider-mean carry-forward on the retained key metrics, but Stage2 already contains an Anthropic-side comparison reversal and the boundary layer stays sparse/non-uniform \\
\bottomrule
\end{tabular}
\caption{NRR-Patterns status ladder under the split comparison contract. Status is assigned from the exact-replication spine plus the separated robustness member, while the rows/run column exposes the unequal paired-row counts behind the five-pattern map. Unequal paired-row volume is not equivalent to unequal independent support thickness.}
\label{tab:status_ladder}
\end{table}

\section{Results: Boundary and Control Checks}
\label{sec:results_boundary}

\subsection{Compression Is Not a Drop-In Substitute}

The compression control stays in the main text because it answers the most natural skepticism against a comparison-led delayed-commitment paper: perhaps the gains come only from shorter prompts. The current aggregated comparison, however, does not support that simpler story for the retained three-pattern subset. Table~\ref{tab:r2_takeaway_integrated} shows that principle-only remains clearly positive for all three tested patterns, whereas compression-only is negative for Hierarchical and Delta and only near-neutral for Lazy.

\begin{table}[t]
\centering
\small
\begin{tabular}{lrrr}
\toprule
\textbf{Pattern} & \textbf{Principle-only} & \textbf{Compression-only} & \textbf{Principle+compression} \\
\midrule
Delta & \(+38.148\%\) & \(-25.623\%\) & \(+3.324\%\) \\
Hierarchical & \(+57.212\%\) & \(-8.224\%\) & \(+36.675\%\) \\
Lazy & \(+38.207\%\) & \(+1.523\%\) & \(+8.589\%\) \\
\bottomrule
\end{tabular}
\caption{Provider-balanced weighted reduction for the retained r2 compression comparison. Values are read from the aggregated r2 comparison summary artifact described in Sec.~\ref{sec:surfaces}.}
\label{tab:r2_takeaway_integrated}
\end{table}

The comparative structure matters. Principle-only remains the strongest read for all three tested patterns, while the combination condition stays positive but weaker than principle-only. The provider-separated control rows also show why this section must remain condition-bounded rather than aggregate-only: for Hierarchical, compression-only is negative on Anthropic and OpenAI but positive on Gemini, while for Lazy it is positive on Anthropic and OpenAI but strongly negative on Gemini. Compression-only therefore functions as a useful control for this retained subset, but not as an equivalent or provider-invariant substitute for delayed-commitment patterning on that tested three-pattern slice.

\begin{figure}[t]
\centering
\includegraphics[width=0.95\textwidth]{../figures/fig_r2_comparison_weighted.png}
\caption{Compression control for the retained three-pattern subset. Principle-only remains strongest, while compression-only weakens or fails for key cases.}
\label{fig:paper7_r2_comparison}
\end{figure}

\subsection{Selected Stage B Boundary Checks}

Selected Stage B evidence is carried into this manuscript only to make the comparison map honest about where provider-separated quality behavior weakens, reverses, or becomes non-effective. This Stage B layer is the selected carry-forward subset of the historical `nrr-boundary' confirmatory suite, not a new manuscript-native experiment block. The retained \texttt{misinterpretation\_rate} plus \texttt{rework\_cost\_tokens} pair is not presented as an original manuscript preregistration freeze; it is a paper-local integration rule that keeps the smallest downstream quality/cost pair needed for honesty reporting. Table~\ref{tab:stageb_selected_integrated} summarizes the two key metrics retained for this role.

\begin{table}[t]
\centering
\small
\begin{tabular}{lrr}
\toprule
\textbf{Pattern} & \textbf{Misinterpretation improvement} & \textbf{Rework improvement (tokens)} \\
\midrule
Conditional & \(+0.0781\) & \(+122.6481\) \\
Hierarchical & \(+0.0120\) & \(-12.1852\) \\
Delta & \(+0.0028\) & \(-26.6296\) \\
Branch & \(-0.0006\) & \(-24.7222\) \\
Lazy & \(-0.1242\) & \(-234.3704\) \\
\bottomrule
\end{tabular}
\caption{Selected Stage B provider-balanced key-metric improvements, read from the Stage B provider-balanced artifact described in Sec.~\ref{sec:surfaces}. Positive values indicate improvement over baseline.}
\label{tab:stageb_selected_integrated}
\end{table}

This boundary-honesty layer changes how the comparison map should be read:
\begin{itemize}
  \item Conditional has the clearest positive provider-mean carry-forward on the retained Stage B key metrics, but that carry-forward remains sparse rather than uniformly active.
  \item Lazy is the clearest degraded pattern on the carried-forward confirmatory surface, with negative provider-balanced means on both retained key metrics.
  \item Hierarchical, Delta, and Branch show provider-sensitive weakening or reversal cells rather than a clean all-provider carry-forward.
  \item Sign-flip and non-effective regions must be reported explicitly rather than hidden behind pooled summaries.
\end{itemize}

The provider-separated read makes the same point more sharply. Conditional remains positive on provider means for misinterpretation and rework across Anthropic, OpenAI, and Gemini, but the carry-forward is still sparse rather than uniformly active: in \path{nrr-boundary/stats/stageb_all/stageb_confirmatory_tests_holm30.csv}, misinterpretation carries \(n_{\mathrm{nonzero}}=25/54\) and rework carries \(33/54\), and the provider-separated rows still contain substantial zero-count mass plus some negative rework cases. By contrast, Hierarchical flips negative on Gemini for both retained key metrics, Delta is negative on Anthropic and Gemini for rework, Branch is negative on Gemini for both retained key metrics, and Lazy is strongly negative on Gemini while only weakly positive or near-neutral elsewhere. On the comparison surface, this same provider-local structure is why Branch should be read as a localized instability case rather than as a uniformly collapsing family. This is why Stage B is kept as a boundary-honesty layer: it does not replace the main comparison spine, but it prevents readers from over-reading the positive comparison map as provider-invariant.

Sign-flip evidence is also direct rather than rhetorical. In the current Stage B boundary artifact, sign-flip boundaries appear in 20 of 45 pattern-metric cells. Conditional shows no Stage B sign-flip cells under this setup, while Branch, Delta, Hierarchical, and Lazy all show sign-flip boundaries on multiple metrics. Figure~\ref{fig:paper7_stageb_signflip} is included here for compact honesty reporting rather than for boundary-paper completeness.

\begin{figure}[t]
\centering
\includegraphics[width=0.95\textwidth]{../figures/paper7_stageb_sign_flip_boundaries.png}
\caption{Selected Stage B sign-flip boundary map carried forward as the honesty layer for NRR-Patterns.}
\label{fig:paper7_stageb_signflip}
\end{figure}

The confirmatory results are consistent with this reading. In \path{nrr-boundary/stats/stageb_all/stageb_confirmatory_tests_holm30.csv}, only five Holm-significant cells survive correction: two positive cells for Conditional (misinterpretation and rework) and three degraded cells for Lazy (update consistency, internal coherence, and self-contradiction). That pattern is exactly the kind of selective carry-forward the manuscript needs: enough to make the positive comparison story honest, but not enough to turn the manuscript into a second confirmatory paper on operator boundaries.

Non-effective reporting is also part of the current manuscript-facing surface rather than a package-only afterthought. Table~\ref{tab:stageb_noneffective_integrated} lists the selected Stage B metrics whose provider-separated artifact shows all-provider zero cells under the current setup. These cells do not overturn the main comparison read by themselves, but they do mark regions where a carry-forward interpretation should stop short of claiming an active effect.

\begin{table}[t]
\centering
\small
\begin{tabular}{L{2.4cm}L{8.8cm}}
\toprule
\textbf{Pattern} & \textbf{All-provider zero metrics on selected Stage B} \\
\midrule
Conditional & internal coherence error rate; trace completeness \\
Delta & internal coherence error rate; trace completeness \\
Branch & internal coherence error rate; trace completeness \\
Hierarchical & none in the selected all-provider-zero slice \\
Lazy & none in the selected all-provider-zero slice \\
\bottomrule
\end{tabular}
\caption{Selected Stage B all-provider non-effective cells. A cell appears here when the provider-separated honesty artifact reports zero-direction behavior for Anthropic, Gemini, and OpenAI under the current setup. Provider-specific zero or near-zero cells remain in the shipped provider-separated artifact even when they do not rise to the all-provider slice shown here.}
\label{tab:stageb_noneffective_integrated}
\end{table}

\section{Discussion}
\label{sec:discussion}

\subsection{What the Pattern Map Supports}

The supported read is condition-bounded utility, not universal superiority. On the exact-replication spine, all five patterns are aggregate-positive relative to their own paired baselines, but they do not all carry forward equally well once the provider-separated stability table, the separated robustness member, and provider-sensitive boundary evidence are added. Hierarchical and Delta remain robustly positive on the split comparison surface. Lazy remains positive on the exact-replication spine and the separated robustness member, but its carried-forward boundary read is degraded rather than robustly positive. Conditional stays aggregate-positive but already contains a Stage2 provider-side reversal. Branch fails even earlier by showing exact-replication Anthropic instability and then losing positive status on the robustness member itself through a tail-risk-dominated provider-local failure. That is why the right manuscript output is a family-benchmark pattern map, not a same-scale ranking.

\subsection{Why Provider Sensitivity Changes the Read}

Provider-balanced summaries are useful for the exact-replication headline, but the separated robustness member and provider-separated views are necessary where aggregation can hide reversals or tail-risk concentration. The integrated manuscript therefore uses four evidence roles on purpose: exact replication for the headline family-benchmark read, provider-separated stability for comparison-side reversals and concentration, robustness for controlled condition-shift stress, and boundary honesty for the carry-forward caveat layer. Collapsing those roles would either overstate the comparison or overinflate the boundary section into a second main paper.

\subsection{Why Compression Is Not Enough}

The compact control against generic prompt compression matters because it distinguishes delayed-commitment patterning from a simpler prompt shortening story. On the retained r2 surface, principle-only stays clearly positive where compression-only weakens or fails for the three tested patterns. This does not prove that compression is never useful; it shows that the tested subset gains are not adequately explained by prompt shortening alone.

\subsection{How to Read Unstable, Sign-Flip, and Non-Effective Regions}

These regions are not embarrassing leftovers. They are part of the contribution because they define where reuse becomes unsafe, unsupported, or too weak to generalize under the current evidence state. In this paper, unstable is primarily a comparison-side label, while sign-flip and non-effective are boundary-honesty labels. Keeping those roles distinct is part of the manuscript's scientific honesty.

\subsection{Series Position}

NRR-Patterns is the pattern-comparison and boundary-honesty layer upstream of Energy. It is not itself the calibration paper, and it is not the downstream closure paper. Its job is to make the upstream pattern read clean enough that later calibration and closure work do not inherit an internally blurred comparison story.

\section{Limitations}
\label{sec:limitations}

This manuscript cites traceable comparison and Stage B values directly from shipped claim-trace artifacts, and the manuscript-facing naming freeze is recorded in the integrated manifests. The release surface is auditable rather than a full rerun bundle: it ships the manuscript assets, the comparison and boundary manifests, a release-surface checksum manifest, and the named comparison/boundary artifacts needed for inspection. In this package, exact replication means the fixed matched-run output surface can be audited and drilled into from shipped summaries plus preserved run-specific members; it does not mean the review ZIP is itself a full independent rerun environment. The exact-replication headline has one package-level summary artifact for inspection while preserving the upstream run-specific members for drill-down, and the separated robustness member remains visible rather than folded into that headline mean. Provider sensitivity remains real and must not be hidden by a single aggregate summary, and the five-pattern map is carried by unequal paired-row volume across patterns. That paired-row volume is useful for auditability, but it is not the same thing as independent support thickness. Ordered composition is intentionally demoted out of the default mainline, and no-position ablation remains package-first rather than main-text essential. These are scope controls rather than unfinished claim-layer work.

\section{Conclusion}
\label{sec:conclusion}

NRR-Patterns supports a condition-bounded pattern map over five delayed-commitment patterns under a split comparison contract. On the exact-replication spine, all five patterns remain aggregate-positive relative to their own paired baselines, but the provider-separated stability summary and the separated temperature-robustness member show that those positives are not equally stable. Hierarchical and Delta remain robustly positive on the split comparison surface, Lazy stays positive on comparison and robustness while degrading on the carried-forward boundary layer, Conditional remains aggregate-positive but already contains a Stage2 provider-side reversal, and Branch becomes the clearest fragile provider-local case because exact-replication Anthropic instability is followed by a Gemini-led, tail-risk-dominated Stage2 failure rather than by uniform collapse across providers. That five-pattern map should be read with explicit paired-row-volume asymmetry because Delta and Conditional carry fewer paired rows per run than Hierarchical, Lazy, and Branch, but paired-row volume is not the same thing as independent support thickness. A compact compression control over the retained three-pattern subset shows that those tested gains are not reducible to generic prompt shortening alone. Selected Stage B carry-forward then makes that map honest by showing that Conditional has the clearest positive provider-mean boundary-facing read, but still with sparse/non-uniform activation, while Lazy degrades and several other patterns contain provider-sensitive reversal or all-provider non-effective cells. The supported output is therefore not a universal ranking, but a comparison-led family benchmark map with explicit provider-separated stability, robustness, sign-flip, and non-effective regions.

\section{Reproducibility and Appendix Home}
\label{sec:repro}

The integrated manuscript release surface has an auditable claim-trace home rather than a full end-to-end rerun bundle. The package-level naming freeze is recorded in \path{repro/paper7_integrated_package_map_v19_2026-04-01.md}, the paper-local split comparison contract is recorded in \path{repro/paper7_split_comparison_contract_v12_2026-04-01.md}, the paper-local Stage B selection rule is recorded in \path{repro/paper7_stageb_selection_rule_v1_2026-03-31.md}, the release-manifest checksum anchor is \path{repro/paper7_integrated_repro_checksums_sha256.txt}, and the shipped audit-surface checksum manifest is \path{repro/paper7_integrated_audit_surface_sha256_v19_2026-04-01.txt}. The main comparison claim-trace members are:
\begin{itemize}
  \item \path{results/analysis/paper7_integrated_comparison_headline_summary_v2_2026-03-31.csv}
  \item \path{results/analysis/paper7_provider_separated_stability_summary_v2_2026-04-01.csv}
  \item \path{results/paper7_results_allproviders_clean.zip::results/analysis/paper7_allproviders_pattern_summary_robust_t00.csv}
  \item \path{results/paper7_results_allproviders_clean.zip::results/analysis/paper7_allproviders_paired_deltas_merged_t00.csv}
  \item \path{results/paper7_results_stage1_t00_rep2.zip::results/analysis/paper7_allproviders_pattern_summary_robust_t00.csv}
  \item \path{results/paper7_results_stage2_t03.zip::results/analysis/paper7_allproviders_pattern_summary_robust_t03.csv}
  \item \path{results/paper7_results_stage2_t03.zip::results/analysis/paper7_allproviders_paired_deltas_merged_t03.csv}
  \item \path{results/paper7_results_stage2_t03.zip::results/analysis/paper7_prod_gemini_clean_t03_paired_deltas.csv}
  \item \path{results/analysis/paper7_comparison_summary_r2_aggregated.csv}
\end{itemize}

The carried-forward boundary-honesty claim-trace members are:
\begin{itemize}
  \item \path{nrr-boundary/stats/stageb_all/mst_provider_balanced.csv}
  \item \path{nrr-boundary/stats/stageb_all/mst_provider_separated.csv}
  \item \path{nrr-boundary/stats/stageb_all/mst_sign_flip_boundaries.csv}
  \item \path{nrr-boundary/stats/stageb_all/stageb_confirmatory_tests_holm30.csv}
\end{itemize}

Appendix/package-first surfaces include the no-position ablation
(\path{results/paper7_no_position_t00.zip}, \path{results/analysis/paper7_no_position_summary.csv}).
An ordered-combination follow-up surface is also retained outside the main-text
claim hierarchy, but it is named for traceability only and is omitted from the
current audit-surface zip/default release surface. These package-first surfaces
are therefore not promoted into the integrated mainline, and only the no-position
ablation is shipped in the current audit package.

The integrated manuscript checksum manifest is
\path{manuscript/checksums_active_review_surface_sha256.txt}. In the current
release surface, \path{scripts/verify_current_package.sh} verifies the manuscript
asset set together with the release manifests, while
\path{scripts/verify_audit_surface.py} verifies the shipped claim-trace surface
against the audit checksum manifest. Together, the package map, the split
comparison contract, the Stage B selection rule, the two claims manifests, the
integrated comparison headline summary, the release-manifest checksum anchor,
the audit-surface checksum manifest, and the integrated manuscript checksum
manifest define the current integrated manuscript package surface for this manuscript.

\section*{Acknowledgments}
The author acknowledges language refinement and drafting support from Claude
(Anthropic), ChatGPT/Codex (OpenAI), and Gemini (Google). All claims and
decisions in this manuscript remain the sole responsibility of the author.

\begin{thebibliography}{99}\sloppy

\bibitem{saito2025nrr}
Saito, K. (2025).
\textit{NRR-Core: Non-Resolution Reasoning as a Computational Framework for Contextual Identity and Ambiguity Preservation}.
arXiv:2512.13478.

\bibitem{saito2026phi}
Saito, K. (2026).
\textit{NRR-Phi: Text-to-State Mapping for Ambiguity Preservation in LLM Inference}.
arXiv:2601.19933.

\bibitem{saito2026ime}
Saito, K. (2026).
\textit{NRR-IME: Interface Decomposition for Stable State Transitions on Stateless LLM APIs}.
Series-path repository tree snapshot cited for context only. Manuscript snapshot in cited tree: \path{manuscript/current/paper4-nrr-ime-v67.tex}. \url{https://github.com/kei-saito-research/nrr-ime/tree/99f5ebbbc3e854b575f95b2e70ef65686a44ba4d}.

\bibitem{saito2026transfer}
Saito, K. (2026).
\textit{Cross-Task Parser/Update Reuse of a Fixed Hidden-State Operator-Target Interface on Stateless LLM APIs}.
Series-path repository tree snapshot cited for context only. Manuscript snapshot in cited tree: \path{manuscript/current/paper5-nrr-transfer-v123.tex}. \url{https://github.com/kei-saito-research/nrr-transfer/tree/24bd24ea183f282cf535205ca4447e8169bac3b0}.

\bibitem{saito2026coupled}
K.~Saito,
\newblock ``NRR-Coupled (cp): Dependency-Consistency Evaluation for Coupled State Updates,''
\newblock Series-path repository tree snapshot cited for context only. Manuscript snapshot in cited tree:
\path{manuscript/current/paper6-nrr-coupled-v21.tex}.
\url{https://github.com/kei-saito-research/nrr-coupled/tree/aab7cc500acc413d4f585c8a50a9045cb5477eb5}.

\bibitem{xu2023recomp}
F.~Xu, W.~Shi, and E.~Choi (2023).
\textit{RECOMP: Improving Retrieval-Augmented LMs with Compression and Selective Augmentation}.
arXiv:2310.04408.

\bibitem{chevalier2023compress}
A.~Chevalier et al. (2023).
\textit{Adapting Language Models to Compress Contexts}.
arXiv:2305.14788.

\bibitem{liang2022helm}
P.~Liang et al. (2023). \textit{Holistic Evaluation of Language Models}.
\emph{Transactions on Machine Learning Research}.

\bibitem{zheng2023mtbench}
L.~Zheng et al.,
\newblock ``Judging LLM-as-a-Judge with MT-Bench and Chatbot Arena,''
\newblock \emph{NeurIPS Datasets and Benchmarks}, 2023.

\bibitem{dubois2024alpacaeval}
Y.~Dubois et al.,
\newblock ``Length-Controlled AlpacaEval: A Simple Way to Debias Automatic Evaluators,''
\newblock \emph{arXiv preprint arXiv:2404.04475}, 2024.

\bibitem{kojima2022zeroshot}
T.~Kojima et al.,
\newblock ``Large Language Models are Zero-Shot Reasoners,''
\newblock \emph{NeurIPS}, 2022.

\bibitem{wang2023selfconsistency}
X.~Wang et al.,
\newblock ``Self-Consistency Improves Chain of Thought Reasoning in Language Models,''
\newblock \emph{ICLR}, 2023.

\bibitem{zhao2021calibrate}
Z.~Zhao et al.,
\newblock ``Calibrate Before Use: Improving Few-Shot Performance of Language Models,''
\newblock \emph{ICML}, 2021.

\end{thebibliography}

\end{document}
